% =====================================================================
% ECIR 2027 short paper (v3): retrieval-utility framing, minimal new runs.
% LLNCS macro package, Version 2.25 of 2026/09/03.
% Short papers: up to 6 pages plus references; double-anonymous review.
%
% HOW TO USE THIS DRAFT (details in the student runbook)
%  \tbd{KEY}          a number still to be produced; KEY is listed in the runbook.
%  \verify{VAL}{KEY}  an existing number to recompute (Holm families changed).
%  \TODO{...}         text to write or decide once a result is in.
%  \branch{S}{T}{F}   outcome-dependent text; set boolean S below using the
%                     runbook's decision rules. Draft builds show both versions.
%  Lines tagged  % DROP-IF-NOT-RUN(<task>)  are deleted if that task does not run.
%  Set \booltrue{submission} for the final build: any placeholder left then
%  stops compilation with an error naming it.
%  Fill numbers with:  python forge_ecir_tools.py fill forge_ecir2027_v3.tex placeholders.json
%
% REGISTRATION ABSTRACT (plain text for EasyChair; holds for any outcome):
%  Retrieval-augmented generation is the default remedy when a large language
%  model must write code in a domain-specific language (DSL) it has seen little
%  of. Retrievers rank exemplars by their relevance to the request; a DSL's
%  toolchain can instead judge the program each exemplar leads to, giving a
%  deterministic, reference-free measure of retrieval utility. We use this
%  signal to study exemplar retrieval for SysML v2, Solidity, and Modelica. On
%  500 requirements per language, paired ablations swap retrievers across
%  languages and replace retrieved exemplars with random in-domain ones. We then
%  build a small test collection in which every pooled exemplar is labeled by
%  the compiler's verdict on the program it produces, evaluate lexical and dense
%  retrievers against these utility labels, and test whether retrieval scores
%  can decide when to retrieve and whether a reranker trained on the labels
%  selects more useful exemplars. We release the labeled pools for
%  utility-oriented evaluation of DSL retrieval.
% =====================================================================
\documentclass[runningheads]{llncs}
%
\usepackage[T1]{fontenc}
\usepackage{newtxtext}
\usepackage[varvw]{newtxmath}
\usepackage{graphicx}
\usepackage{booktabs}
\usepackage{multirow}
\usepackage{xspace}
\usepackage{xcolor}
\usepackage{etoolbox}
\usepackage{hyperref}
\renewcommand\UrlFont{\color{blue}\rmfamily}
\urlstyle{rm}

\newcommand{\pp}{\,pp\xspace}

% ---------------- placeholder machinery ----------------
\newbool{submission}
\boolfalse{submission}   % \booltrue{submission} for the final build

% Outcome switches: set from the runbook's decision rules.
\newbool{Rone} \boolfalse{Rone} % R1: true if Solidity A1 with BM25+lineage diversity is NOT significantly below A0
\newbool{Rtwo} \boolfalse{Rtwo} % R2: true if relevance-ranked A1 significantly beats random-exemplar A1 in >= 2 of 3 languages
\newbool{Urel} \boolfalse{Urel} % U1: true if the best lexical retriever's Gain@5 significantly exceeds random exemplars' in both Solidity and Modelica
\newbool{Sel}  \boolfalse{Sel}  % A2: true if the cross-validated selective policy significantly beats the better fixed policy in >= 1 language
\newbool{Rrk}  \boolfalse{Rrk}  % R7: true if R7 ran and the reranker significantly beats the deployed retriever in >= 1 language

\newcommand{\tbd}[1]{\ifbool{submission}%
  {\PackageError{forge}{Unfilled placeholder #1}{See the runbook placeholder table.}}%
  {\mbox{\textcolor{red}{\textsf{\tiny[#1]}}}}}
\newcommand{\verify}[2]{\ifbool{submission}%
  {\PackageError{forge}{Unverified value #2}{See the runbook.}}%
  {\textcolor{orange}{#1}}}
\newcommand{\TODO}[1]{\ifbool{submission}%
  {\PackageError{forge}{Unresolved TODO}{See the runbook.}}%
  {\textcolor{red}{[\textbf{TODO:} #1]}}}
\newcommand{\branch}[3]{\ifbool{submission}%
  {\ifbool{#1}{#2}{#3}}%
  {\mbox{\textcolor{blue}{\textsf{\tiny[#1 true:]}}}~#2~\mbox{\textcolor{blue}{\textsf{\tiny[#1 false:]}}}~#3}}
% --------------------------------------------------------

\begin{document}
\emergencystretch=3em
%
\title{When Do Retrieved Exemplars Help Domain-Specific Code Generation? Compiler Outcomes as Retrieval Utility}
\titlerunning{Compiler Outcomes as Retrieval Utility}
%
% Double-anonymous submission: no author names, affiliations, or identifying links.
\author{Anonymous Author(s)}
\authorrunning{Anonymous}
\institute{Anonymous Institution}
%
\maketitle
%
\begin{abstract}
Retrieval-augmented generation (RAG) is the default remedy when a large language model must write code in a domain-specific language (DSL) it has seen little of. Retrievers rank exemplars by relevance to the request; a DSL's toolchain can judge the program each exemplar leads to, giving a deterministic, reference-free measure of retrieval utility. We use this signal to study exemplar retrieval for SysML~v2, Solidity, and Modelica. On 500 requirements per language, lexical retrieval of natural-language--code exemplars raises Modelica compile validity by 13.7\pp and lowers Solidity's by 11.4\pp.
\branch{Rone}{Swapping retrievers between the two languages moves Solidity's change to $\tbd{R1.sol.d}$\pp, so the sign follows retriever design.}{The Solidity loss persists under Modelica's retriever ($\tbd{R1.sol.d}$\pp), so the sign follows the language.}
\branch{Rtwo}{Relevance-ranked exemplars beat random in-domain ones by at least \tbd{R2.mingap}\pp.}{Random in-domain exemplars come within \tbd{R2.maxgap}\pp of relevance-ranked ones.}
Labeling \tbd{U1.nlabels} pooled exemplars by the compiler's verdict on the program each produces alone,
\branch{Urel}{we find that lexical relevance predicts utility, yet \tbd{U1.sol.hurt.lex}\% of top-ranked Solidity exemplars make the program worse than no retrieval.}{we find that top-ranked exemplars are barely more useful than random ones, and \tbd{U1.sol.hurt.lex}\% of top-ranked Solidity exemplars make the program worse than no retrieval.}
\branch{Sel}{Deciding whether to retrieve from the top-1 retrieval score recovers \tbd{A2.recov}\% of the per-requirement oracle gain on held-out requirements.}{The top-1 retrieval score is too weak a predictor to decide when to retrieve.}
\branch{Rrk}{A reranker trained on the compiler labels raises validity by $\tbd{R7.best.d}$\pp over the deployed retriever.}{}
We release the labeled pools as a test collection for utility-oriented DSL retrieval.

\keywords{Retrieval-augmented generation \and Retrieval utility \and Code generation \and Domain-specific languages.}
\end{abstract}
%
\section{Introduction}
\label{sec:intro}

LLMs generate general-purpose code far more reliably than code in domain-specific languages (DSLs)~\cite{chen2021evaluating,joel2024lowresource}, which are underrepresented in pretraining data and whose correctness is decided by a toolchain. The standard IR remedy is RAG~\cite{lewis2020rag}: condition generation on retrieved documentation or natural-language (NL)--code exemplars~\cite{zhou2023docprompting,dutta2024rar}. Retrievers rank this context by relevance to the request, while the user needs a valid program, and the two can diverge: misaligned context degrades code generation~\cite{gu2025alliancecoder,wan2025modelica}, and random documents can raise RAG accuracy~\cite{cuconasu2024noise}. Measuring utility directly usually needs reference outputs or an LLM judge~\cite{salemi2024erag}.

DSLs come with a better judge. A compiler, a test generator, or a simulator decides for every generated program whether it is valid, deterministically and without references. We use this judge to study exemplar retrieval in SysML~v2, Solidity, and Modelica, asking:
\textbf{RQ1:} Does lexical exemplar retrieval improve one-shot DSL generation, and does its effect follow the retriever, the language, or relevance?
\textbf{RQ2:} How well does relevance, lexical or dense, predict an exemplar's utility as judged by the toolchain?
\textbf{RQ3:} Can toolchain outcomes decide when to retrieve and which exemplars to retrieve?
We contribute paired retrieval ablations across three DSLs; a test collection of \tbd{U1.nlabels} pooled exemplars for 200 requirements, each labeled by the toolchain's verdict on the program it produces; and evaluations of selective retrieval and of a reranker trained on these labels.

\section{Related Work}
\label{sec:related}

\textbf{Retrieval for code generation.} Retrieved documentation~\cite{zhou2023docprompting} or code exemplars~\cite{dutta2024rar} help with unfamiliar libraries and low-resource languages~\cite{cassano2024knowledge}, with gains that depend on the retrieved source~\cite{wang2025coderagbench}; CEDAR selects code demonstrations by embedding or lexical similarity~\cite{nashid2023cedar}. DSL generators target one language each~\cite{cibrian2025agent,bouamra2025systemp,peng2025soleval,luo2025fsm,xiang2025modigen,wan2025modelica}, and compilable models can still be semantically wrong~\cite{hendricks2025text2model}.

\textbf{Relevance and utility.} eRAG labels each retrieved document by the downstream score an LLM reaches with that document alone~\cite{salemi2024erag}; our labels follow this protocol with a toolchain as the metric. Retrievers of in-context examples can be trained from LM feedback, such as the likelihood of a target output~\cite{rubin2022epr}, distilled LLM preferences~\cite{wang2024llmr}, or a black-box LM's scores~\cite{shi2024replug}; these signals need reference outputs or model likelihoods.

\textbf{When to retrieve.} Adaptive retrieval decides per query whether to retrieve at all~\cite{mallen2023adaptive,asai2024selfrag,jiang2023flare}, a query performance prediction problem~\cite{cronen2002qpp}. We test whether a retrieval score supports that decision when the toolchain defines success.

\section{Experimental Setup}
\label{sec:setup}

\textbf{Tasks and corpora.} SysML~v2: 1{,}543 requirements written by Gemini~2.5~Pro from Wikipedia-listed engineered systems, with 300 curated exemplar models. Solidity: 1{,}500 requirements in 21 application domains, expanded by Gemini~2.5~Flash from DefiLlama protocol entries, with 1{,}500 human-authored exemplar contracts (libraries, audit contests, on-chain code) described by Gemini. Modelica: 1{,}560 robotics tasks (13 robot families) with 1{,}500 separately authored exemplars. We use a fixed sample of 500 requirements per language (498 for Modelica). Requirements and exemplars come from disjoint sources: \tbd{A10.leak.proto} DefiLlama protocols appear among exemplar sources, and \tbd{A10.leak.dup} generated contracts have a near-duplicate exemplar (token 5-gram Jaccard $\geq 0.8$).

\textbf{Generation.} Every condition is one GLM-5.2 call through OpenRouter at temperature 0.2 with pinned providers (\tbd{D7.providers}). A0 prompts with the requirement alone; A1 adds $k=5$ retrieved exemplars, their code truncated to 80 non-comment lines, and for SysML~v2 three passages from a 1{,}775-chunk index of the language specification. Modelica requirements are first normalized into an intermediate representation and an execution contract in both conditions.

\textbf{Retrievers.} \emph{Binary cosine}: cosine similarity of binary token sets between the requirement and exemplar descriptions (deployed for SysML~v2 and Solidity). \emph{BM25}~\cite{robertson2009bm25} over description, category, and tags, keeping at most one exemplar per source lineage (deployed for Modelica). \emph{Dense}: \tbd{D9.dense} embeddings without fine-tuning, used for pooling and offline evaluation. \emph{Random}: exemplars drawn uniformly from the same corpus under the same truncation (on SysML~v2 also three random specification chunks).

\textbf{Utility judge.} Each generated program receives a graded gain: 0 if it fails to compile; 1 if it compiles; 2 if it also passes interface-level checks, namely fuzzing and boundary tests generated programmatically from the contract ABI (Solidity; single-file \texttt{solc}~\tbd{D2.solc}) or a successful FMU build and simulation (Modelica). Set-level results report compile validity.

\textbf{Utility-labeled pools.} For 100 requirements each in Solidity and Modelica, drawn at random from the 500, the pool is the union of the three retrievers' top-5 plus two random exemplars (\tbd{U1.poolsize} exemplars per requirement on average). Each pooled exemplar is placed alone ($k=1$) in the A1 prompt, and the resulting program's gain is its label. A0 is sampled three times per pooled requirement; an exemplar \emph{helps} or \emph{hurts} when its gain is above or below the median A0 gain. Relabeling a random 10\% of pairs with a second sample reproduces \tbd{U2.agree}\% of gains (Cohen's $\kappa=\tbd{U2.kappa}$). Retrievers are scored by Gain@5 (mean gain of their top five), Help@5, Hurt@5, and nDCG@5 against the pooled gains.

\textbf{Statistics.} Set-level comparisons are paired per requirement: exact McNemar tests, Holm-corrected within each table, with paired bootstrap 95\% confidence intervals (10{,}000 resamples). Pool comparisons use paired randomization tests over requirements (10{,}000 sign flips). Code, requirement samples, labeled pools, and outputs: \TODO{anonymized repository URL}.

\section{Results}
\label{sec:results}

\begin{table}[t]
\centering
\caption{Compile validity (\%) of one-shot generation, paired within each column ($n=500$, 500, 498). The deployed retriever is binary cosine for SysML~v2 and Solidity and BM25 with lineage diversity for Modelica; the alternative is the other one. Selective: retrieve only when the top-1 score clears a threshold chosen by 10-fold cross-validation. Oracle: the better of A0 and A1 per requirement.}
\label{tab:set}
\scriptsize
\setlength{\tabcolsep}{6pt}
\begin{tabular}{@{}lrrr@{}}
\toprule
Condition & SysML~v2 & Solidity & Modelica \\
\midrule
A0: no retrieval & \tbd{R4.sys.a0} & 69.6 & 49.2 \\
A1: deployed retriever & 15.2 & 58.2 & 62.9 \\
A1: alternative retriever & -- & \tbd{R1.sol.cv} & \tbd{R1b.mod.cv} \\
A1: random exemplars & \tbd{R2.sys.cv} & \tbd{R2.sol.cv} & \tbd{R2.mod.cv} \\
\midrule
Selective (deployed retriever) & \tbd{A2.sys.sel} & \tbd{A2.sol.sel} & \tbd{A2.mod.sel} \\
Oracle, max(A0, A1) & \tbd{A2.sys.orc} & \tbd{A2.sol.orc} & \tbd{A2.mod.orc} \\
\bottomrule
\end{tabular}
\end{table}

\begin{table}[t]
\centering
\caption{Solidity compiler errors per sample on the first attempt, by error class (paired, $n=500$).}
\label{tab:err}
\scriptsize
\setlength{\tabcolsep}{6pt}
\begin{tabular}{@{}lrrrr@{}}
\toprule
Error class & A0 & A1 deployed & A1 alternative & A1 random \\
\midrule
Unresolved import & \tbd{A1.imp.a0} & \tbd{A1.imp.a1} & \tbd{A1.imp.alt} & \tbd{A1.imp.rand} \\
Compiler version or pragma & \tbd{A1.ver.a0} & \tbd{A1.ver.a1} & \tbd{A1.ver.alt} & \tbd{A1.ver.rand} \\
Undeclared identifier & \tbd{A1.und.a0} & \tbd{A1.und.a1} & \tbd{A1.und.alt} & \tbd{A1.und.rand} \\
Type error & \tbd{A1.type.a0} & \tbd{A1.type.a1} & \tbd{A1.type.alt} & \tbd{A1.type.rand} \\
Parser error & \tbd{A1.parse.a0} & \tbd{A1.parse.a1} & \tbd{A1.parse.alt} & \tbd{A1.parse.rand} \\
Other & \tbd{A1.other.a0} & \tbd{A1.other.a1} & \tbd{A1.other.alt} & \tbd{A1.other.rand} \\
\midrule
Total & 0.82 & 1.38 & \tbd{R1.sol.err} & \tbd{R2.sol.err} \\
\bottomrule
\end{tabular}
\end{table}

\textbf{RQ1: Set-level effect of retrieval} (Table~\ref{tab:set}). Deployed retrieval raises Modelica compile validity by 13.7\pp (Holm $p=\verify{2.0\times10^{-6}}{V.p.mod}$) and lowers Solidity's by 11.4\pp ($p=\verify{1.4\times10^{-3}}{V.p.sol}$), raising Solidity compiler errors per sample from 0.82 to 1.38, mostly through \tbd{A1.top.class} (\tbd{A1.top.share}\% of the increase; Table~\ref{tab:err}) \TODO{A1: mechanism in one clause, and whether the alternative and random retrievers show the same classes}. On SysML~v2, exemplars and specification passages raise validity from \tbd{R4.sys.a0}\% to 15.2\% ($p=\tbd{R4.sys.p}$).
\branch{Rone}{With Modelica's BM25 lineage-diversity retriever, Solidity's change becomes $\tbd{R1.sol.d}$\pp ($p=\tbd{R1.sol.p}$), so the sign of retrieval's effect follows the retriever.}{With Modelica's BM25 lineage-diversity retriever, Solidity's change is still $\tbd{R1.sol.d}$\pp ($p=\tbd{R1.sol.p}$), so the sign follows the language.}
Binary cosine on Modelica gives \tbd{R1b.mod.cv}\%.
\branch{Rtwo}{Random in-domain exemplars fall \tbd{R2.mingap}\pp or more below relevance-ranked ones, so relevance matters beyond in-domain context.}{Random in-domain exemplars come within \tbd{R2.maxgap}\pp of relevance-ranked ones, so in-domain context of any kind carries most of the effect.}

\begin{table}[t]
\centering
\caption{Exemplar utility on the labeled pools (100 requirements per language). Each exemplar is used alone; Gain@5 is the mean graded gain (0--2) of a method's top five, Help@5 and Hurt@5 the share of them scoring above or below the median A0 program, and nDCG@5 is computed against the pooled gains. Random: the two random exemplars per requirement. Reranker: 5-fold cross-validation by requirement.}
\label{tab:pool}
\scriptsize
\setlength{\tabcolsep}{3.5pt}
\begin{tabular}{@{}lrrrrrrrr@{}}
\toprule
 & \multicolumn{4}{c}{Solidity (A0 gain \tbd{U1.sol.a0.gain})} & \multicolumn{4}{c}{Modelica (A0 gain \tbd{U1.mod.a0.gain})} \\
\cmidrule(lr){2-5}\cmidrule(l){6-9}
Method & Gain@5 & Help@5 & Hurt@5 & nDCG@5 & Gain@5 & Help@5 & Hurt@5 & nDCG@5 \\
\midrule
Binary cosine & \tbd{U1.sol.cos.gain} & \tbd{U1.sol.cos.help} & \tbd{U1.sol.cos.hurt} & \tbd{U1.sol.cos.ndcg} & \tbd{U1.mod.cos.gain} & \tbd{U1.mod.cos.help} & \tbd{U1.mod.cos.hurt} & \tbd{U1.mod.cos.ndcg} \\
BM25 + lineage diversity & \tbd{U1.sol.bm25.gain} & \tbd{U1.sol.bm25.help} & \tbd{U1.sol.bm25.hurt} & \tbd{U1.sol.bm25.ndcg} & \tbd{U1.mod.bm25.gain} & \tbd{U1.mod.bm25.help} & \tbd{U1.mod.bm25.hurt} & \tbd{U1.mod.bm25.ndcg} \\
Dense & \tbd{U1.sol.dense.gain} & \tbd{U1.sol.dense.help} & \tbd{U1.sol.dense.hurt} & \tbd{U1.sol.dense.ndcg} & \tbd{U1.mod.dense.gain} & \tbd{U1.mod.dense.help} & \tbd{U1.mod.dense.hurt} & \tbd{U1.mod.dense.ndcg} \\
Random & \tbd{U1.sol.rand.gain} & \tbd{U1.sol.rand.help} & \tbd{U1.sol.rand.hurt} & -- & \tbd{U1.mod.rand.gain} & \tbd{U1.mod.rand.help} & \tbd{U1.mod.rand.hurt} & -- \\
\midrule
Utility-trained reranker & \tbd{A13.sol.rr.gain} & \tbd{A13.sol.rr.help} & \tbd{A13.sol.rr.hurt} & \tbd{A13.sol.rr.ndcg} & \tbd{A13.mod.rr.gain} & \tbd{A13.mod.rr.help} & \tbd{A13.mod.rr.hurt} & \tbd{A13.mod.rr.ndcg} \\
\bottomrule
\end{tabular}
\end{table}

\textbf{RQ2: Relevance and exemplar utility} (Table~\ref{tab:pool}).
\branch{Urel}{Relevance carries utility signal: the best lexical retriever's top five score $\tbd{U1.sol.lexgap}$ and $\tbd{U1.mod.lexgap}$ gain above random exemplars on Solidity and Modelica ($p\le\tbd{U1.p}$), yet \tbd{U1.sol.hurt.lex}\% of its top five Solidity exemplars still produce a worse program than no retrieval.}{Relevance carries little utility signal: the best lexical retriever's top five differ from random exemplars by $\tbd{U1.sol.lexgap}$ and $\tbd{U1.mod.lexgap}$ gain on Solidity and Modelica, and \tbd{U1.sol.hurt.lex}\% of its top five Solidity exemplars produce a worse program than no retrieval.}
Dense retrieval reaches nDCG@5 of \tbd{U1.sol.dense.ndcg} and \tbd{U1.mod.dense.ndcg} \TODO{U1: compare with the lexical retrievers, citing the randomization p}. Across pooled exemplars, harm is best predicted by \tbd{A13.harm.feat} (odds ratio \tbd{A13.harm.or}, 95\% CI $[\tbd{A13.harm.orci}]$), the same mechanism behind the set-level Solidity loss \TODO{A13: one concrete harmful exemplar, e.g.\ a top-ranked contract whose library import the generated code copies, failing single-file compilation}. The set level agrees with the pools where they overlap: \TODO{U1/R1: the retriever with the higher Gain@5 also gives the higher A1 validity in Table~\ref{tab:set}, or say where they disagree}.

\begin{figure}[t]
\centering
\IfFileExists{fig-selective.pdf}%
  {\includegraphics[width=0.7\textwidth]{fig-selective.pdf}}%
  {\ifbool{submission}{\PackageError{forge}{Missing figure fig-selective.pdf}{Run the curve command of the helper script.}}%
    {\fbox{\parbox[c][3.4cm][c]{0.66\textwidth}{\centering\textcolor{red}{Figure placeholder: \texttt{forge\_ecir\_tools.py curve fig-selective.pdf \dots}}}}}}
\caption{Selective retrieval: compile validity when only the requirements with the highest top-1 retrieval scores retrieve, from none (A0, left) to all (A1, right); dotted lines mark the per-requirement oracle. The curve is in-sample; Table~\ref{tab:set} reports the cross-validated policy.}
\label{fig:sel}
\end{figure}

\textbf{RQ3: When and what to retrieve.} Figure~\ref{fig:sel} shows where retrieval pays off \TODO{A2: describe the curve in one sentence, e.g.\ validity peaks when the top X\% retrieve}.
\branch{Sel}{Retrieving only when the top-1 score clears a cross-validated threshold beats the better fixed policy by $\tbd{A2.sel.d}$\pp on \tbd{A2.sel.dsl} (95\% CI $[\tbd{A2.sel.ci}]$, $p=\tbd{A2.sel.p}$) and recovers \tbd{A2.recov}\% of the gap to the per-requirement oracle (Table~\ref{tab:set}).}{Retrieving only when the top-1 score clears a cross-validated threshold never beats the better fixed policy (Table~\ref{tab:set}), although the per-requirement oracle lies \tbd{A2.orcgap}\pp above it: the top-1 score is a weak predictor of when retrieval helps.}
A reranker trained on the pooled labels (logistic regression over the three retrieval scores and exemplar features) reaches cross-validated nDCG@5 of \tbd{A13.sol.rr.ndcg} and \tbd{A13.mod.rr.ndcg} \TODO{A13: against the best single retriever, with p}.
Trained on all 100 labeled requirements and applied to the held-out requirements (400 Solidity, 398 Modelica), its top five give \tbd{R7.sol.cv}\% and \tbd{R7.mod.cv}\% compile validity against \tbd{R7.sol.base}\% and \tbd{R7.mod.base}\% for the deployed retriever ($p=\tbd{R7.p}$), \branch{Rrk}{so compiler-derived labels transfer to generation.}{so the offline gains did not transfer to generation.} % DROP-IF-NOT-RUN(R7)

\section{Discussion and Scope}
\label{sec:discussion}

\textbf{Retrieval under repair.} In a four-model harness built on these prompts, synthesis and at most two compiler-repair rounds raise validity from the A1 level to 34.2\%, 92.0\%, and 81.3\% (SysML~v2, Solidity, Modelica), so repair absorbs much of retrieval's harm on Solidity. Within that configuration the odds ratio of final validity per standard deviation of top-1 score is \tbd{A14.sol.or} ($[\tbd{A14.sol.orci}]$) on Solidity and \tbd{A14.mod.or} ($[\tbd{A14.mod.orci}]$) on Modelica; this association is observational, and fewer repair rounds are a natural further target for utility-trained retrieval.

\textbf{Implications for IR.} Toolchain verdicts give graded utility judgments for the cost of one generation and one compilation per requirement--exemplar pair, with no assessors, reference outputs, or LLM judges. Pools like ours evaluate any DSL retriever offline, and the same labels can train retrievers the way LM likelihoods do~\cite{rubin2022epr,wang2024llmr,shi2024replug}, with an executable signal in their place.

\textbf{Scope.} We use one backbone, one sample per requirement--exemplar pair, 100 labeled requirements per language, and untuned retrievers. Labels score each exemplar alone, which ignores interactions among the five exemplars of a prompt; Table~\ref{tab:set} is the set-level check. Compilation and interface-level checks are necessary conditions for correctness; they do not establish that a program implements its requirement.
\TODO{If R7 did not run: replace the R7 sentence in RQ3 and the abstract with the offline result only.}

\section{Conclusion}
\label{sec:conclusion}

A DSL toolchain turns every generation into a utility judgment, which lets retrieval be evaluated by what it does to the program. Used this way, lexical exemplar retrieval helped on Modelica and hurt on Solidity, with a sign set by \branch{Rone}{the retriever.}{the language.}
\branch{Rtwo}{Relevance-ranked exemplars beat random in-domain ones,}{Random in-domain exemplars did nearly as well as relevance-ranked ones,}
and relevance predicted the utility of individual exemplars \branch{Urel}{only in part.}{poorly.}
\branch{Rrk}{A reranker trained on compiler labels improved generation, closing the loop from judgment to retrieval.}{Training retrievers on compiler labels at larger scale is the natural next step.}

\begin{credits}
\subsubsection{\discintname}
The authors have no competing interests to declare that are relevant to the content of this article.
\end{credits}
%
% ---- Bibliography ----
% references.bib = the v1 bibliography; references_v2_additions.bib = entries added since.
\bibliographystyle{splncs04}
\bibliography{references,references_v2_additions}
%
\end{document}
